% Môn: Vật lý
% Lớp: 10
% Chương: Chương I. Mở đầu
% Bài: L10C1 Bài 2. Các quy tắc an toàn trong phòng thực hành Vật lí · Nhận biết nội quy và quy tắc an toàn
% Số câu: 185
% App đọc file này trực tiếp — sửa rồi Commit trên GitHub, bấm Đồng bộ GitHub .tex

% L10C1 Bài 2. Các quy tắc an toàn trong phòng thực hành Vật lí



% ===== Câu 84 =====
% ID: L10C1B2-19-DS
% Mức: NB
\dangbt{Nhận biết nội quy và quy tắc an toàn}
\begin{ex}
% ID: L10C1B2-19-DS
% Mức: NB
Trong tình huống trước khi bắt đầu thí nghiệm đối với nội quy chung, hãy đánh giá các phát biểu sau.
\choiceTF
{\True Cần đọc hướng dẫn, kiểm tra dụng cụ và chỉ thao tác khi được giáo viên cho phép.}
{Có thể tự ý thay đổi bố trí thí nghiệm nếu thao tác diễn ra nhanh.}
{\True Phải báo cho giáo viên hoặc người phụ trách khi phát hiện nguy cơ vượt khả năng kiểm soát.}
{\True Chỉ tiếp tục thí nghiệm sau khi nguyên nhân nguy hiểm đã được loại bỏ hoặc kiểm soát.}
\loigiai{
\begin{itemchoice}
\item (Đúng) Cần đọc hướng dẫn, kiểm tra dụng cụ và chỉ thao tác khi được giáo viên cho phép. (Vì): vì đây là biện pháp an toàn của dạng nội quy chung; b sai vì hành vi “tự ý thay đổi bố trí thí nghiệm” vẫn nguy hiểm; c và d đúng do sự cố phải được báo cáo, cô lập và kiểm soát trước khi tiếp tục
\item (Sai) Có thể tự ý thay đổi bố trí thí nghiệm nếu thao tác diễn ra nhanh.
\item (Đúng) Phải báo cho giáo viên hoặc người phụ trách khi phát hiện nguy cơ vượt khả năng kiểm soát.
\item (Đúng) Chỉ tiếp tục thí nghiệm sau khi nguyên nhân nguy hiểm đã được loại bỏ hoặc kiểm soát.
\end{itemchoice}
}
\end{ex}

% ===== Câu 86 =====
% ID: L10C1B2-20-DS
% Mức: NB
\begin{ex}
% ID: L10C1B2-20-DS
% Mức: NB
Trong tình huống khi phát hiện thiết bị có dấu hiệu bất thường đối với nội quy chung, hãy đánh giá các phát biểu sau.
\choiceTF
{\True Cần đọc hướng dẫn, kiểm tra dụng cụ và chỉ thao tác khi được giáo viên cho phép.}
{Có thể tự ý thay đổi bố trí thí nghiệm nếu thao tác diễn ra nhanh.}
{\True Phải báo cho giáo viên hoặc người phụ trách khi phát hiện nguy cơ vượt khả năng kiểm soát.}
{\True Chỉ tiếp tục thí nghiệm sau khi nguyên nhân nguy hiểm đã được loại bỏ hoặc kiểm soát.}
\loigiai{
\begin{itemchoice}
\item (Đúng) Cần đọc hướng dẫn, kiểm tra dụng cụ và chỉ thao tác khi được giáo viên cho phép. (Vì): vì đây là biện pháp an toàn của dạng nội quy chung; b sai vì hành vi “tự ý thay đổi bố trí thí nghiệm” vẫn nguy hiểm; c và d đúng do sự cố phải được báo cáo, cô lập và kiểm soát trước khi tiếp tục
\item (Sai) Có thể tự ý thay đổi bố trí thí nghiệm nếu thao tác diễn ra nhanh.
\item (Đúng) Phải báo cho giáo viên hoặc người phụ trách khi phát hiện nguy cơ vượt khả năng kiểm soát.
\item (Đúng) Chỉ tiếp tục thí nghiệm sau khi nguyên nhân nguy hiểm đã được loại bỏ hoặc kiểm soát.
\end{itemchoice}
}
\end{ex}

% ===== Câu 88 =====
% ID: L10C1B2-21-DS
% Mức: TH
\begin{ex}
% ID: L10C1B2-21-DS
% Mức: TH
Trong tình huống trong lúc thay đổi cách bố trí dụng cụ đối với nội quy chung, hãy đánh giá các phát biểu sau.
\choiceTF
{\True Cần đọc hướng dẫn, kiểm tra dụng cụ và chỉ thao tác khi được giáo viên cho phép.}
{Có thể tự ý thay đổi bố trí thí nghiệm nếu thao tác diễn ra nhanh.}
{\True Phải báo cho giáo viên hoặc người phụ trách khi phát hiện nguy cơ vượt khả năng kiểm soát.}
{\True Chỉ tiếp tục thí nghiệm sau khi nguyên nhân nguy hiểm đã được loại bỏ hoặc kiểm soát.}
\loigiai{
\begin{itemchoice}
\item (Đúng) Cần đọc hướng dẫn, kiểm tra dụng cụ và chỉ thao tác khi được giáo viên cho phép. (Vì): vì đây là biện pháp an toàn của dạng nội quy chung; b sai vì hành vi “tự ý thay đổi bố trí thí nghiệm” vẫn nguy hiểm; c và d đúng do sự cố phải được báo cáo, cô lập và kiểm soát trước khi tiếp tục
\item (Sai) Có thể tự ý thay đổi bố trí thí nghiệm nếu thao tác diễn ra nhanh.
\item (Đúng) Phải báo cho giáo viên hoặc người phụ trách khi phát hiện nguy cơ vượt khả năng kiểm soát.
\item (Đúng) Chỉ tiếp tục thí nghiệm sau khi nguyên nhân nguy hiểm đã được loại bỏ hoặc kiểm soát.
\end{itemchoice}
}
\end{ex}

% ===== Câu 89 =====
% ID: L10C1B2-21-TL
% Mức: VD
\begin{bt}
% ID: L10C1B2-21-TL
% Mức: VD
Xây dựng một bảng kiểm ngắn gồm ít nhất bốn mục cho tình huống trong lúc thay đổi cách bố trí dụng cụ liên quan đến nội quy chung.
\loigiai{
Bảng kiểm gồm: đọc hướng dẫn; nhận diện mối nguy; kiểm tra dụng cụ và bảo hộ; thực hiện đọc hướng dẫn, kiểm tra dụng cụ và chỉ thao tác khi được giáo viên cho phép; theo dõi bất thường; thu dọn và báo cáo sau thí nghiệm.
}
\end{bt}

% ===== Câu 90 =====
% ID: L10C1B2-22-DS
% Mức: TH
\begin{ex}
% ID: L10C1B2-22-DS
% Mức: TH
Trong tình huống sau khi hoàn thành phép đo đối với nội quy chung, hãy đánh giá các phát biểu sau.
\choiceTF
{\True Cần đọc hướng dẫn, kiểm tra dụng cụ và chỉ thao tác khi được giáo viên cho phép.}
{Có thể tự ý thay đổi bố trí thí nghiệm nếu thao tác diễn ra nhanh.}
{\True Phải báo cho giáo viên hoặc người phụ trách khi phát hiện nguy cơ vượt khả năng kiểm soát.}
{\True Chỉ tiếp tục thí nghiệm sau khi nguyên nhân nguy hiểm đã được loại bỏ hoặc kiểm soát.}
\loigiai{
\begin{itemchoice}
\item (Đúng) Cần đọc hướng dẫn, kiểm tra dụng cụ và chỉ thao tác khi được giáo viên cho phép. (Vì): vì đây là biện pháp an toàn của dạng nội quy chung; b sai vì hành vi “tự ý thay đổi bố trí thí nghiệm” vẫn nguy hiểm; c và d đúng do sự cố phải được báo cáo, cô lập và kiểm soát trước khi tiếp tục
\item (Sai) Có thể tự ý thay đổi bố trí thí nghiệm nếu thao tác diễn ra nhanh.
\item (Đúng) Phải báo cho giáo viên hoặc người phụ trách khi phát hiện nguy cơ vượt khả năng kiểm soát.
\item (Đúng) Chỉ tiếp tục thí nghiệm sau khi nguyên nhân nguy hiểm đã được loại bỏ hoặc kiểm soát.
\end{itemchoice}
}
\end{ex}

% ===== Câu 91 =====
% ID: L10C1B2-22-TL
% Mức: VD
\begin{bt}
% ID: L10C1B2-22-TL
% Mức: VD
Một học sinh đã tự ý thay đổi bố trí thí nghiệm. Hãy nêu trình tự xử lí ngay sau khi phát hiện hành vi này.
\loigiai{
Trình tự: yêu cầu dừng ngay; không tiếp cận nếu chưa an toàn; cô lập nguồn nguy hiểm; báo giáo viên; sơ cứu hoặc xử lí theo hướng dẫn; chỉ hoạt động lại sau khi được xác nhận an toàn.
}
\end{bt}

% ===== Câu 93 =====
% ID: L10C1B2-23-TL
% Mức: VD
\begin{bt}
% ID: L10C1B2-23-TL
% Mức: VD
So sánh biện pháp phòng ngừa và biện pháp ứng phó đối với nguy cơ thuộc nhóm nội quy chung.
\loigiai{
Phòng ngừa diễn ra trước sự cố, tập trung vào kiểm tra, loại bỏ hoặc giảm nguy cơ. Ứng phó diễn ra khi sự cố đã hoặc sắp xảy ra, gồm dừng hoạt động, cô lập nguồn nguy hiểm, báo cáo và sơ cứu. Hai nhóm biện pháp bổ sung cho nhau.
}
\end{bt}

% ===== Câu 95 =====
% ID: L10C1B2-24-TL
% Mức: VDC
\begin{bt}
% ID: L10C1B2-24-TL
% Mức: VDC
Đề xuất cách tổ chức nhóm học sinh để thực hiện nhiệm vụ về nội quy chung an toàn và có thể kiểm tra được.
\loigiai{
Phân công trưởng nhóm kiểm tra điều kiện, người thao tác, người quan sát an toàn và người ghi chép. Dùng bảng kiểm, xác nhận trước mỗi bước, quy định tín hiệu dừng và báo ngay bất thường.
}
\end{bt}

% ===== Câu 139 =====
% ID: L10C1B2-46-TN
% Mức: NB
\begin{ex}
% ID: L10C1B2-46-TN
% Mức: NB
Khi phát hiện thiết bị có dấu hiệu bất thường liên quan đến nội quy chung, hành động nào an toàn nhất?
\choice
{tự ý thay đổi bố trí thí nghiệm}
{tiếp tục thao tác rồi báo cáo sau}
{nhờ bạn khác làm thay mà không kiểm tra điều kiện an toàn}
{\True đọc hướng dẫn, kiểm tra dụng cụ và chỉ thao tác khi được giáo viên cho phép}
\loigiai{
Hành động đúng là “đọc hướng dẫn, kiểm tra dụng cụ và chỉ thao tác khi được giáo viên cho phép” vì giúp phòng ngừa sai thao tác và phát hiện nguy cơ trước khi thực hành. Chọn D.
}
\end{ex}

% ===== Câu 143 =====
% ID: L10C1B2-48-TN
% Mức: TH
\begin{ex}
% ID: L10C1B2-48-TN
% Mức: TH
Sau khi hoàn thành phép đo liên quan đến nội quy chung, hành động nào an toàn nhất?
\choice
{nhờ bạn khác làm thay mà không kiểm tra điều kiện an toàn}
{\True đọc hướng dẫn, kiểm tra dụng cụ và chỉ thao tác khi được giáo viên cho phép}
{tự ý thay đổi bố trí thí nghiệm}
{tiếp tục thao tác rồi báo cáo sau}
\loigiai{
Hành động đúng là “đọc hướng dẫn, kiểm tra dụng cụ và chỉ thao tác khi được giáo viên cho phép” vì giúp phòng ngừa sai thao tác và phát hiện nguy cơ trước khi thực hành. Chọn B.
}
\end{ex}

% ===== Câu 145 =====
% ID: L10C1B2-49-TN
% Mức: TH
\begin{ex}
% ID: L10C1B2-49-TN
% Mức: TH
Khi một bạn chưa đọc hướng dẫn liên quan đến nội quy chung, hành động nào an toàn nhất?
\choice
{\True đọc hướng dẫn, kiểm tra dụng cụ và chỉ thao tác khi được giáo viên cho phép}
{tự ý thay đổi bố trí thí nghiệm}
{tiếp tục thao tác rồi báo cáo sau}
{nhờ bạn khác làm thay mà không kiểm tra điều kiện an toàn}
\loigiai{
Hành động đúng là “đọc hướng dẫn, kiểm tra dụng cụ và chỉ thao tác khi được giáo viên cho phép” vì giúp phòng ngừa sai thao tác và phát hiện nguy cơ trước khi thực hành. Chọn A.
}
\end{ex}

% ===== Câu 146 =====
% ID: L10C1B2-50-TN
% Mức: TH
\begin{ex}
% ID: L10C1B2-50-TN
% Mức: TH
Khi khu vực làm việc bị ướt liên quan đến nội quy chung, hành động nào an toàn nhất?
\choice
{tự ý thay đổi bố trí thí nghiệm}
{tiếp tục thao tác rồi báo cáo sau}
{nhờ bạn khác làm thay mà không kiểm tra điều kiện an toàn}
{\True đọc hướng dẫn, kiểm tra dụng cụ và chỉ thao tác khi được giáo viên cho phép}
\loigiai{
Hành động đúng là “đọc hướng dẫn, kiểm tra dụng cụ và chỉ thao tác khi được giáo viên cho phép” vì giúp phòng ngừa sai thao tác và phát hiện nguy cơ trước khi thực hành. Chọn D.
}
\end{ex}

% ===== Câu 148 =====
% ID: L10C1B2-51-TN
% Mức: TH
\begin{ex}
% ID: L10C1B2-51-TN
% Mức: TH
Khi cần rời bàn thí nghiệm liên quan đến nội quy chung, hành động nào an toàn nhất?
\choice
{tiếp tục thao tác rồi báo cáo sau}
{nhờ bạn khác làm thay mà không kiểm tra điều kiện an toàn}
{\True đọc hướng dẫn, kiểm tra dụng cụ và chỉ thao tác khi được giáo viên cho phép}
{tự ý thay đổi bố trí thí nghiệm}
\loigiai{
Hành động đúng là “đọc hướng dẫn, kiểm tra dụng cụ và chỉ thao tác khi được giáo viên cho phép” vì giúp phòng ngừa sai thao tác và phát hiện nguy cơ trước khi thực hành. Chọn C.
}
\end{ex}

% ===== Câu 150 =====
% ID: L10C1B2-53-TN
% Mức: VD
\begin{ex}
% ID: L10C1B2-53-TN
% Mức: VD
Khi kết quả đo khác dự kiến liên quan đến nội quy chung, hành động nào an toàn nhất?
\choice
{\True đọc hướng dẫn, kiểm tra dụng cụ và chỉ thao tác khi được giáo viên cho phép}
{tự ý thay đổi bố trí thí nghiệm}
{tiếp tục thao tác rồi báo cáo sau}
{nhờ bạn khác làm thay mà không kiểm tra điều kiện an toàn}
\loigiai{
Hành động đúng là “đọc hướng dẫn, kiểm tra dụng cụ và chỉ thao tác khi được giáo viên cho phép” vì giúp phòng ngừa sai thao tác và phát hiện nguy cơ trước khi thực hành. Chọn A.
}
\end{ex}

% ===== Câu 151 =====
% ID: L10C1B2-54-TN
% Mức: VD
\begin{ex}
% ID: L10C1B2-54-TN
% Mức: VD
Khi giáo viên yêu cầu dừng thao tác liên quan đến nội quy chung, hành động nào an toàn nhất?
\choice
{tự ý thay đổi bố trí thí nghiệm}
{tiếp tục thao tác rồi báo cáo sau}
{nhờ bạn khác làm thay mà không kiểm tra điều kiện an toàn}
{\True đọc hướng dẫn, kiểm tra dụng cụ và chỉ thao tác khi được giáo viên cho phép}
\loigiai{
Hành động đúng là “đọc hướng dẫn, kiểm tra dụng cụ và chỉ thao tác khi được giáo viên cho phép” vì giúp phòng ngừa sai thao tác và phát hiện nguy cơ trước khi thực hành. Chọn D.
}
\end{ex}

